\documentclass[12pt]{article}
\usepackage[utf8]{inputenc}
\usepackage[T1]{fontenc}
\usepackage{hyperref}
\title{Narragansett Language Resources for K-12 Curriculum Development}
\author{Compiled by Claude Sonnet 5.5 (claude-sonnet-5-5) for Andr\'as Kornai}
\date{}
\begin{document}
\maketitle

\section{Introduction}
Narragansett (ISO 639-3 xnt) is an Algic (Eastern Algonquian) language of what is now Rhode Island, United States. It is dormant, with no remaining first-language speakers. The tribe has begun revival work based on early books and manuscripts, according to Wikipedia. This search did not identify a specific school programme or modern dictionary.

\section{Primary Educational Resources}
The search found no modern primers, textbooks or school programmes. Roger Williams's 1643 book is the main historical teaching-style source, although Williams described it as an ``implicit dialogue'' rather than a dictionary.

\section{Linguistic and Reference Materials}
\begin{itemize}
\item Narragansett language, Wikipedia overview: \url{https://en.wikipedia.org/wiki/Narragansett_language}
\item Roger Williams, \emph{A Key Into the Language of America} (1643), Project Gutenberg: \url{https://m.gutenberg.org/ebooks/63701}
\item \emph{A key into the language of America}, Internet Archive scan: \url{https://archive.org/details/bub_gb_wOfpAPRxlVYC}
\item Roger Williams, English-to-Algonquian collection, American Antiquarian Society: \url{https://collections.americanantiquarian.org/EnglishtoAlgonquian/rogerwilliams}
\end{itemize}

\section{Religious and Community Resources}
No scripture or community-organisation resources were found in these searches.

\section{Digital and Audio Resources}
No audio, video or software resources were found in these searches. As a dormant language, Narragansett has no native-speaker recordings.

\section{Curriculum Development Context}
Revival rests on 17th-century and later written sources, mainly Williams's \emph{Key}. The search did not locate the tribe's own language programme or any modern teaching materials. A follow-up search on the Narragansett Indian Tribe's language programme is needed before curriculum planning.

\end{document}